\section[算术局部对称空间上的 Hecke 代数]{算术局部对称空间上的 Hecke 代数（4 月 6 日）}

\srcpage{53}

上次我们讨论了 Hecke 代数：若 $G$ 作用在空间 $X$ 上、$K\subset G$，商空间 $X/K$ 没有 $K$-作用，但 $X/K$ 上的函数有 $K\backslash G/K$ 上函数的作用，这是一个群胚代数。

我们特别关心 $\Bun_G(C)$ 或其在数域上的类比这样的例子。我们把它们呈现为双陪集空间 $\Gamma\backslash G_\R/K$，例如 $\mathrm{SL}_2(\Z)\backslash \mathrm{SL}_2(\R)/\mathrm{SO}_2$。让我们把它纳入 Hecke 故事。

此情形 $X=\Gamma\backslash G_\R$，所以 Hecke 代数应是 $K\backslash G_\R/K$ 上的函数。例如取 $\mathrm{SL}_2(\Z)\backslash\mathrm{SL}_2(\R)/\mathrm{SO}_2$，双陪集空间 $\mathrm{SO}_2\backslash\mathrm{SL}_2(\R)/\mathrm{SO}_2$ 是什么？首先 $\mathrm{SL}_2(\R)/\mathrm{SO}_2$ 是上半平面，$\mathrm{SO}_2$ 从左侧绕 $i$ 旋转：左 $\mathrm{SO}_2$ 是 $i\in\mathbb{H}$ 在 $\mathrm{SL}_2(\R)$ 中的稳定化子。商同胚于 $\R_+$，即 $\mathbb{H}$ 中从 $i$ 出发的所有（双曲度量下的）距离。

\srcpage{54}

Hecke 对称告诉我们，从双陪集得到 $G_\R/K$ 上函数的算子，它们与 $G_\R$-作用交换。固定 $r\in\R_+$；在函数 $f:\mathbb{H}\to\C$ 上得到的算子把 $f$ 在半径 $r$ 的圆上平均：

\[(\cO_r*f)(z)=\oint_{|w-z|=r} f(w).\]

这些算子彼此交换，且与 $\mathrm{SL}_2(\R)$-作用交换。函数的正则性不是重点，但我们要积分收敛，所以 $L^2$ 可用。特别地，因为它们与 $\mathrm{SL}_2(\R)$-作用交换，它们下降到商 $\Gamma\backslash\mathbb{H}$，其中 $\Gamma\subset\mathrm{SL}_2(\R)$ 是离散子群。

另一表现是

\[K\backslash G/K=G\backslash(G/K\times G/K),\]

其中 $G$ 对角作用。这引向微分算子——空间 $X$ 上的微分算子是函数上的局部自同态。这里的局部是强意义：若 $\cO$ 是微分算子，$(\cO*f)(x)$ 只依赖 $f$ 在 $x$ 的 Taylor 级数。这与之前相反：固定距离的平均非常非局部。这个微分算子定义归功于 Grothendieck；它是在代数几何中说清微分算子的好办法。从核变换的观点看，微分算子的积分核支撑在 $X\times X$ 中对角线的无穷小邻域内。

\begin{theorem}[Harish-Chandra]
设 $G_\R$ 是实 Lie 群，$K$ 是其极大紧子群。则 $G_\R/K$ 上 $G_\R$-等变微分算子代数是交换环，事实上是多项式代数 $\C[\mathfrak{a}]^W$。
\end{theorem}

其中 $\mathfrak{a}$ 是 $G_\R/K$ 的 Cartan Lie 代数，$W$ 是小 Weyl 群。

\begin{example}
对 $\mathbb{H}=\mathrm{SL}_2(\R)/\mathrm{SO}_2$，微分算子代数是 $\C[\Delta_{\mathbb{H}}]$，其中 $\Delta_{\mathbb{H}}$ 是双曲度量的 Laplacian。
\end{example}

高维有更多算子，但它们都交换。

这个微分算子代数是 Hecke 代数，此时通过微分作用；它由核支撑在对角线邻域内这一点保证。于是 $\Gamma\backslash G_\R/K$ 上带有自然的交换对称代数 $\mathrm{Diff}_{G_\R}(G_\R/K)$。

在取曲线上所有点的限制乘积的情形

\[G(F)\backslash \prod_{x\in C}{}^{\prime} G(K_x)\Big/ G(\oo_x),\]

我们对每个 $x$ 有「局部」Hecke 代数作用，它们交换，通过张量所有这些局部作用给出巨大的代数作用。局部 Hecke 代数联系于双陪集空间 $G(\oo_x)\backslash G(K_x)/G(\oo_x)$。

\begin{theorem}[Satake]
Hecke 代数 $C_c[G(\oo_x)\backslash G(K_x)/G(\oo_x)]$ 也是多项式代数。
\end{theorem}

数域情形有 $K_x=\Q_p$、$\oo_x=\Z_p$ 之类，考察 $G(\Z_p)\backslash G(\Q_p)/G(\Z_p)$ 随 $p$ 变化。函数域情形 $K_x=\F_q((t))$、$\oo_x=\F_q[[t]]$，即 Laurent 级数对 Taylor 级数取双商。

\begin{example}
$G=\mathrm{PGL}_2$ 时这些 Hecke 算子长什么样？此情形双商曾是 $\mathbb{H}$ 中半径的空间；局部域情形的图景对数域与函数域竟然无所谓。

固定 $K=\Q_p$；$K=\F_p((t))$ 的情形会惊人地类似。则 $G(K)/G(\oo)=\mathrm{PGL}_2(\Q_p)/\mathrm{PGL}_2(\Z_p)$，它与 $\Q_p^2$ 中 $\Z_p$-格模去位似构成的空间双射（前几讲见过类似的整体例子）。确切地说，这样的格是 $\Z_p$-子模 $\Lambda\subset\Q_p^2$，使得当 $p$ 可逆时得到整个 $\Q_p^2$：即

\[\Lambda\otimes_{\Z_p}\Q_p\cong \Q_p^2.\]

任何格都可被 $\mathrm{GL}_2(\Q_p)$ 的作用带到标准格；「位似」指把格在重新缩放下等同。

格的例子有哪些？一个是 $\Lambda_0:=\Z_p^2\subset\Q_p^2$，由标准基 $\{e_1,e_2\}$ 张成。也可以取 $p\Lambda_0:=\mathrm{span}\{pe_1,pe_2\}$。或者取 $\{pe_1,e_2\}$ 的张成。
\end{example}

\srcpage{55}

我们画一个图，顶点是 $\Q_p^2$ 中的格模去位似（即陪集空间的点），若在模去位似后存在包含

\[\Lambda\subset\Lambda'\subset p\Lambda,\]

就把 $\Lambda$ 与 $\Lambda'$ 连一条边。

这是一个等价关系——为得到对称性，例如把 $\Lambda'$ 重新缩放为 $p^{-1}\Lambda'$。

所得的图称为 $\mathrm{PGL}_2$ 在 $\Q_p$ 上的 building（建筑）\footnote{「building」是对称空间概念的 $p$ 进类比。}；事实上它是一棵树。与 $\Lambda_0$ 相邻的格等价于 $\Lambda_0/p\Lambda_0\cong\F_p^2$ 中的直线，也等价于 $\mathbb{P}^1(\F_p)$ 中的点。所以 $\Lambda_0$ 有 $p+1$ 个邻居。

对任意 $\Lambda$ 重复此分析，可断定这个图是无限 $3$-正则树。原文此处配图 1（$\mathrm{SL}_2$ 的 Bruhat–Tits 树，来源为 nLab 的 Casselman 讲义，并附「TODO：补引用」），译本从略。\yizhu{原文图 1 为 $\mathrm{SL}_2$ 的 Bruhat–Tits 树示意图，来源为 nLab 上的 Casselman 讲义；此处未复刻图片，图注中的 TODO 已并入正文。}{A（原文 p.55）}

这棵树很像单位圆盘的 $p$ 进类比，它就是数域情形中双曲空间的类比物。顶点 $G(K)/G(\oo)$ 是称为 $G$ 的仿射 Grassmann 簇的空间，而双陪集空间 $G(\oo)\backslash G(K)/G(\oo)$ 是 building 上 $G(\oo)$-轨道的空间。我们的情形 $G(\oo)=\mathrm{PGL}_2(\Z_p)$。这些轨道类似于 $\mathbb{H}$ 的情形：两个格在双商中等同当且仅当它们到 $\Lambda_0$ 的距离相同，所以双陪集空间是树度量下到 $\Lambda_0$ 的距离空间。

生成等变微分算子代数的算子是「树 Laplacian」$T_{p^n}$——谱图论中图 Laplacian 的特例。若 $f$ 是顶点集 $G(K)/G(\oo)$ 上的函数，

\[T_{p^n}f(v)=\sum_{d(w,v)=n} f(w).\]

微分算子代数是 $\C[T_p]$。

\srcpage{56}

现在整体化：曲线上每个点都有一个这样的代数，它们不仅交换，而且彼此交换。这个极其巨大的交换代数作用是 Langlands 纲领如此有力的部分原因。

不同点的局部 Hecke 代数彼此交换并不是最令人惊讶的事；真正的奇迹是每一个单独都交换。一旦有了这个，就应取其谱，看看能学到什么。接下来几周我们会花在这两个想法上。

\begin{remark}
定理 19.1 与通常的 Harish-Chandra 同构有何关系？设 $G$ 是实约化群；则定理 19.1 告诉我们 $\mathrm{Diff}_{G\times G}(G)$ 是多项式环。若只关心左不变微分算子，会得到泛包络代数 $\cU(\mathfrak{g})$；右 $G$-作用下降为 $\cU(\mathfrak{g})$ 上的 $G$-作用，所以

\[\mathrm{Diff}_{G\times G}(G)=(\cU(\mathfrak{g}))^G,\]

这其实是 $\cU(\mathfrak{g})$ 的中心。通常的 Harish-Chandra 同构识别

\[Z(\cU(\mathfrak{g}))\cong \C[\mathfrak{h}^*]^W,\]

其中 $\mathfrak{h}$ 是 Cartan 子代数，$W$ 是 Weyl 群。Weyl 群总是反射群，而 Chevalley 定理说反射群下的不变量总是多项式代数。特别地，得到代数 $\C[p_1,\dots,p_\ell]$，其中 $\ell:=\dim\mathfrak{h}$，称为 $G$ 的秩。

Knop [Kno94] 把 Harish-Chandra 同构升级到被 $G$ 作用的球簇（spherical variety）$X$ 的情形。此时 $X$ 上 $G$-等变 $\cD$-模的代数等同于 $\C[\mathfrak{a}_X]^{W_X}$。
\end{remark}

Hecke 代数并不总在一切情形交换——这里发生了某种好事。我们将用场论给出这些代数应当交换的概念性理由。场论中的对象常带有大而有趣的算子代数。

我们谈谈一维量子场论，也叫量子力学。此时有态 Hilbert 空间 $\cH$ 和可观测量的结合代数 $A$，即 $\cH$ 上所有算子。我们设想算子在特定时刻作用。

\begin{center}
\begin{tikzpicture}[scale=1]
\draw[->] (0,0)--(5,0);
\fill (1,0) circle (1.7pt) node[below] {$\psi_0$};
\fill (2.8,0) circle (1.7pt);
\node[above] at (2.8,0.08) {$A$};
\fill (4.3,0) circle (1.7pt) node[below] {$\psi_1$};
\end{tikzpicture}
\end{center}

在这幅图中，算子 $A$ 作用于 $\psi_0$，产生不同的态 $\psi_1$；时间向右流动。

在 Atiyah–Segal 风格的场论中，这由配边编码。若 $V$ 是流形 $M$ 上拓扑场论的态空间，$V'$ 是 $V$ 的线性对偶，则圆柱 $M\times[0,1]$ 解释为配边 $M\sqcup M\to\varnothing$，给出映射

\[e:V\otimes V'\longrightarrow Z(\varnothing)=\C,\]

称为求值（evaluation）。同样，圆柱是配边 $\varnothing\to M\sqcup M'$，给出余求值映射 $c:\C\to V\otimes V'$。这两个映射满足称为 Zorro 之印（mark of Zorro）的关系（TODO：图）。

抽象地说，对称幺半范畴中带有求值与余求值映射、并满足 Zorro 之印关系的对象称为可对偶的；我们已经说明 $V$ 可对偶且 $V'$ 是其对偶。这意味着 $V,V'$ 是有限维向量空间；此时求值是通常的求值映射 $V\otimes V^*\to\C$，$v,\ell\mapsto \ell(v)$，而余求值是

\[c:1\longmapsto \sum_j e_j\otimes e^j,\]

其中 $\{e_j\}$ 是 $V$ 的基，$\{e^j\}$ 是对偶基。余求值映射不依赖基的选择；正是它让有限维要求变得真正清楚。

TODO：算子来自配边（把「双帽」的边界画成一维配边的图……）。我漏掉了这段，抱歉。

现在谈谈界面。若 $V,W$ 是有限维向量空间，可以构造一维 TFT $Z_V,Z_W$，其态空间（在一点上的取值）分别是 $V,W$。（TODO：我看到记号 $Z(\to)=V$、$Z(\to)=W$，可能造成混淆。）TFT 是函子，所以它们之间最天真的态射概念是自然变换，但这不太有用。反之，我们会发现\srcpage{57}考虑类似 Morita 范畴的态射更有用：代数 $A,B$ 之间的态射是 $(B,A)$-双模。

我们想定义「二部场论」：定义在带有（可能为空的）分离余维一子流形的流形上，且一侧染红、另一侧染蓝。物理上，把理论 $Z_V$ 放在红色分量、$Z_W$ 放在蓝色分量，而它们之间的余维一子流形上有某种「界面」数据，这个数据是 TFT 的态射。

TODO：图。

可以证明，二部场论（「从 $V$ 到 $W$ 的界面」）等价于线性映射 $V\to W$。所以带界面的一维 TFT 等价于向量空间。例如 $Z_V$ 与自身的界面是 $\End(V)$ 的元素——可以想成在区间上放一个缺陷或奇点，这个缺陷由 $\End(V)$ 中的算子标记。

从这个观点看，态即元素 $\psi\in V$，它等价于映射 $\C\to V$，可以想成从平凡 TFT 到 $Z_V$ 的界面。我们可以擦掉二部流形的红色部分，并把这数据想成允许我们从带边紧 $1$-流形中提取数，其边界点由 $\psi$ 标记。这有点傻，但当我们升到维度 $2$、再到维度 $4$ 时就不那么傻了：Hecke 正是从这类图景中自然产生的。
